\section{Height conventions and local map inputs}\label{sec:inputs}

The proof combines three inputs: the critical height coefficient from
\cite{X}, the local integration maps of \cite{B}, and the identical-nest
comparison of \cite{S}.  We first fix their common
normalization.  We then specify the analytic spaces in which the local
maps operate. These inputs include finite-graph inequalities, scaling
limits, and local activity maps. The entry construction for the critical
XY model is supplied below.

\subsection{Dual heights and finite comparisons}

For a graph \(G\) consisting of retained nearest-neighbor edges, choose
one orientation of each edge and write \(\partial v\) for the oriented
gradient.  Its unit-conductance Laplacian is \(A_G=\partial^*\partial\),
so that
\begin{equation}\label{inp:energy-convention}
 \langle v,A_Gv\rangle=\sum_{\{x,y\}\in E(G)}(v_y-v_x)^2.
\end{equation}
Thus each undirected edge contributes once.  Heights take values in
\(2\pi\Z\), and an edge difference \(2\pi k\) has weight
\[
 p_b(k)=e^{-b}I_k(b),\qquad
 \sum_{k\in\Z}p_b(k)e^{tk}=\exp\{b(\cosh t-1)\}.
\]
An integral connection \(l\) is an edge cochain with values in
\(2\pi\Z\); it replaces \(\partial h\) by \(\partial h+l\) in
the weights.  Write \(Z_G^l\) for the resulting sum and
\(F_G^l=Z_G^l/Z_G^0\).  Without observations or absolute pins, a
constant translation by \(2\pi\Z\) is divided out on each free
component.  If observations confine the constants, their translation
copies are summed instead.  Changing \(l\) by an integral gradient is
an exact change of height variables.

\begin{lemma}[Fixed-boundary duality]\label{inp:duality}
Work at \(b=b_c\).  Let \(G_M\) be the free graph of the \(M^2\) unit square faces in
\([-n,n]^2\), where \(M=2n\).  If \(l\) has circulation \(2\pi\)
around the central primal vertex and zero circulation around all
other interior primal vertices, then
\[
                         a_n=F_{G_{2n}}^l>0.
\]
The graph retains the edges crossing primal bonds incident to the
fixed spin boundary.
\end{lemma}

\begin{proof}
This is the fixed-spin duality of \cite[Lemma 2.1]{S}.  In its Fourier
current expansion, integration imposes current divergence only at
interior spin vertices.  Bonds lying along the boundary have both
spins fixed and cancel from the ratio.  Every other retained primal
bond crosses one dual bond, including a bond incident to the boundary.
Subtract a unit current from the center to the boundary and rotate the
remaining divergence-free current.  It is a height gradient, uniquely
modulo one common translation, and the subtracted current gives \(l\).
The factors \(e^b\) cancel.  All remaining weights are positive.
\end{proof}

We shall use the following finite-graph rules in exactly this
normalization.  A centered precision means an additional nonnegative
quadratic form, including a limit enforcing centered linear equalities.
A reference cost for a linear statistic
\(X=\sum_e w_e\partial_eh\) is \(\mathcal R(w)=\sum_e w_e^2\).

\begin{proposition}[Finite comparison inputs]\label{inp:comparisons}
The comparisons in \cite[Lemmas 2.2 and
2.3]{X} hold on every fixed primary graph,
also with independent Gaussian observation noises.
\begin{enumerate}
\item Adding centered precision decreases centered covariance and
centered real moment generating functions.  Covariance in an affine
sector, including after a real tilt, dominates centered covariance.
In particular, for a centered inversion-symmetric law,
\(\log\E e^X\ge\frac12\operatorname{Var}(X)\).
\item A shifted-to-centered partition ratio is at most one and
decreases when the same positive precision is added in the two
sectors.  The shifts may be integral connections, prescribed equality
values, or observation centers.  Monotonicity also holds with a fixed
real linear statistic in the numerator exponent; that tilted ratio
need not be at most one.  Infinite penalties may impose nonempty
affine constraints.
\item Invariant centered characteristic coefficients are nonnegative
and increase with precision.  In the finite-graph situations used
here they are positive, and the Hessian of their logarithm at any
phase is at least its Hessian at zero.  A phase on a free component
is invariant whenever its coefficient total is an integer.
\item With any extra centered precision, for each fixed integer
\(r\ge1\) and real \(t\),
\begin{equation}\label{inp:reference-cost}
 \operatorname{Var}(X)\le C_b\mathcal R(w),\qquad
 \E|X|^r\le C_{r,b}\mathcal R(w)^{r/2},\qquad
 \E e^{tX/\sqrt{\mathcal R(w)}}\le C_{t,b}.
\end{equation}
The last assertion is used for \(\mathcal R(w)>0\); a zero-cost
statistic vanishes.  One may minimize the cost over flows with the
prescribed divergence and grounded vertices.
\end{enumerate}
\end{proposition}

The finite Gaussian antecedent is the lattice inequality of Regev and
Stephens-Davidowitz \cite{RegevStephensDavidowitz2017}.  The required
Bessel-weight comparisons are transferred in \cite{X} through
independent discrete Gaussian chains in series: replace a primary edge by \(N\) increments
with normalized weights \((b/(2N))^{k^2}\), impose endpoint and period
constraints, and then let \(N\to\infty\) at fixed graph.  The physical
sum along a chain tends to the law assigning mass \(p_b(k)\) to
\(2\pi k\).  Quadratic observations remain positive
forms in the enlarged coordinates.  This procedure also transfers
the parity identities for two Gaussian replicas and the theta
comparison used for complex observations.  For example, the latter
bounds the imaginary-to-real partition ratio at a real affine shift
by the corresponding centered imaginary ratio;
see \cite[Equation (9.29)]{B} and
\cite[Equation (5.26)]{X}.  The uniform bounds used
here are the primary costs in \eqref{inp:reference-cost}, not inverses
of the microscopic chain precision.

\subsection{The critical Gaussian coefficient}

For the free height square \(B_m=\{0,\ldots,m-1\}^2\), put
\(g_m(x)=x_1/m\).  The structural coefficient is the limit
\[
 a(b)=\lim_{m\to\infty}
       \operatorname{Var}_{B_m,b}\langle h,A_{B_m}g_m\rangle.
\]
The existence and the following identifications are inputs from
\cite{X}.

\begin{proposition}[Height and pin inputs]\label{inp:height}
For the mass threshold \(b_c\) in Theorem~\ref{thm:main},
\(0<b_c<\infty\), \(m(b_c)=0\), and the primary height coefficient
of \cite{X} satisfies
\begin{equation}\label{inp:critical-coefficient}
                              a(b_c)=8\pi.
\end{equation}
At a general positive coefficient \(a=a(b)>0\), the following
statements hold.
\begin{enumerate}
\item On free rectangles tending after scaling to a fixed
nondegenerate rectangle \(U\), the tests
\(\langle h,A_{U_N}f(\cdot/N)\rangle\), for smooth \(f\) near
\(\overline U\), have joint Gaussian Laplace and moment limits with
covariance \(a\int_U\nabla f\cdot\nabla g\), where the lattice
mesh is \(N^{-1}\).  Mean-free density smears have
covariance \(a(-\Delta_{U,N})^{-1}\); bounded piecewise continuous
densities and face smears, centered by a volume smear, are allowed.
\item In a fixed polygonal geometry made of finitely many rectangular
pieces, with positive-width passages and box pin regions \(p\), the
corresponding limits have precision
\[
 a^{-1}\int_U|\nabla v|^2,
 \qquad v\in H^1(U),\quad v=0\ \hbox{on }p.
\]
Free components are taken modulo constants.  A fixed finite-rank
nonnegative quadratic penalty in box averages gives the corresponding
Gaussian penalty; if it controls all free constants, those constants
are integrated in the limit.  Bounded-step coherent rounding is
permitted.
\item The fractional free mean modulo \(2\pi\) tends to uniform
measure, independently of the Gaussian observations.  If \(u_m\)
is the uniform probability on an \(m\)-square and \(t_m\) is a
zero-total profile of bounded reference cost, then for every fixed
nonzero integer \(k\),
\[
                  \E e^{i\langle h,ku_m+t_m\rangle}\longrightarrow0.
\]
Writing \(\bar\rho_m\) for the average of the four face
probabilities and \(z_m=\E e^{i\langle h,\bar\rho_m\rangle}\), one has
\begin{equation}\label{inp:face-ratio}
 \liminf_{m\to\infty}\frac{z_{Gm}}{z_m}
       \ge c_bG^{2-a/(4\pi)}
\end{equation}
for sufficiently large fixed integers \(G\), with \(c_b>0\)
independent of \(G\).
\end{enumerate}
\end{proposition}

These statements are \cite[Sections 1.1 and 6,
Theorems 2.4, 2.5, and 3.1]{X}.
Here and below all mesh limits precede variation of the fixed geometry.
We also use the separated-pin Theorem
3.2 of \cite{X} in the following precise form.
In a centered high law \(\mu_m\), with penalty
\(\frac12\norm{B_mh}^2\) in finitely many component-local box
averages controlling the constants, let \(H,p\) be positively
separated box regions.  Then
\begin{align*}
 D_m(H,p)&=\log\frac{\mu_m(h_H=0\mid h_p=0)}{\mu_m(h_H=0)},\\
 r_{H,m}(b^p)&=
 \frac{\mu_m(h_H=0\mid h_p=b^p)}{\mu_m(h_H=0\mid h_p=0)}
\end{align*}
satisfy \(D_m\ge0\), \(0<r_{H,m}\le1\), and
\(\E_{\mu_m}r_{H,m}=e^{-D_m}\).  The quantities have Gaussian
limits with precision \(a^{-1}\int|\nabla v|^2+\norm{Bv}^2\).
For functions of random guard data, this means approximation in
probability by a sufficiently long fixed list of smooth guard probes,
whose joint law and function converge to the Gaussian counterparts.
The probe list is enlarged after the mesh limit.  Conditional Laplace
transforms and weak laws have this approximation as well, including
when the data are sampled under a fixed finite real linear tilt.
The expectation identity just stated refers to the centered law.

Section 3.6 of \cite{X} supplies the restoration bounds used in
the entry argument.  At zero guard data, a centered restoration ratio
lies between its value in a free patch and its value with a separating
guard pinned.  If all restored bonds and forms are supported in \(H\),
pinning every primary height in \(H\) makes restoration a fixed scalar.
The free-to-guard multiplicative gap is bounded by
\[
                  \frac{\mu_m(h_p=0\mid h_H=0)}{\mu_m(h_p=0)}.
\]
No observation may cross the separating guard in this comparison.

\subsection{Comparing center observables at different sizes}

The information from \cite{S} that we need is weaker than the
asymptotic formula to be proved: it permits an arbitrary slowly varying
factor.  Its sequential form will eventually pass from geometric
square sizes to all integers.

\begin{lemma}[Center comparison]\label{inp:regular-variation}
There is \(c>0\) such that \(a_n\ge cn^{-1/2}\) for all \(n\ge1\).
If \(n_k,n'_k\to\infty\) are integer sequences and
\(n'_k/n_k\to d\in(0,\infty)\), then
\begin{equation}\label{inp:center-ratio}
                        \frac{a_{n'_k}}{a_{n_k}}\longrightarrow d^{-1/8}.
\end{equation}
\end{lemma}

\begin{proof}
Use the identical nests of \cite[Section 6.2]{S}.  Fix a rectilinear
approximation to the disk with relative error \(\varepsilon\), an
annular ratio \(\lambda>1\), and an initial lattice radius \(R_0\).
Let \(R_j=R_0\lambda^j\).  In both squares place exactly the same
centered lattice nest, including its filled core, through the largest
\(N\) with \(R_N\le\eta\min(n_k,n'_k)\).
The bare nest ratio \(F_{\mathrm{nest}}(R_N)\) is identical in the
two observables.

Every subsequence has a further subsequence on which
\(R_N/n_k\to t>0\); then \(R_N/n'_k\to t/d\).
At the fixed parameters both radii stay away from zero and are at
most \(\eta\).  The exterior formula, Equation (13) of \cite{S}, gives, with
\(D=(-1,1)^2\) and
\(G_D(z,w)=(2\pi)^{-1}\log|z-w|^{-1}+R_D(z,w)\),
\[
 \log F_{\mathrm{ext}}(\rho)
 =-\frac{2\pi\log(1/\rho)+4\pi^2R_D(0,0)}{2a}
       +O(\varepsilon+\eta)+o(1).
\]
Its constants are uniform for the limiting radii in question.
The proof of the exterior-attachment Lemma
6.3 applies to each of these two
exteriors: the centered attachment loss depends only on the final
annulus, and the cut and restored exterior energies have the same
regular Green part.  Their error is \(O(\varepsilon+\eta)\),
independently of the fixed sufficiently large annular ratio.
Since \(N\to\infty\), the propagated last-guard density bound also
applies.  Consequently, for any desired accuracy, the logarithm of
each center observable is the sum of its bare nest logarithm and
exterior logarithm, within that accuracy and an \(o(1)\) error.
Choose \(\varepsilon,\eta\) small, then \(\lambda\) large for
the pin interactions, and finally \(R_0\) large for the fixed annuli.
Subtracting cancels the nest and the regular Green term and gives
\[
 \log a_{n'_k}-\log a_{n_k}
             =-\frac\pi a\log d+o(1)
\]
up to an arbitrarily small error.  As \(a=8\pi\), this is
\eqref{inp:center-ratio}; the subsequence argument proves the full
sequential assertion.

For the lower bound, Equations (18)--(19) of \cite{S} give, at fixed
parameters,
\[
 a_n\ge c_1F_{\mathrm{nest}}(R_N),\qquad
 F_{\mathrm{nest}}(R_N)\ge c_2\lambda^{-N/2}
                                  F_{\mathrm{nest}}(R_0).
\]
The core ratio is positive and \(\lambda^N\le Cn\).  This proves
the bound for large \(n\), and positivity handles the remaining sizes.
\end{proof}

\subsection{Covariances, blocks, and analytic norms}

We now fix the reference \(\alpha=\sqrt{8\pi}\).  Its mean-free
Gaussian field \(\psi\) has covariance \(A^+\), the inverse of
the primary unit-edge Laplacian on nonconstant modes; its common mean is independently
uniform modulo \(\omega=2\pi/\alpha\).  On the plane we use only
finite-support covariance formulas.  For a dyadic integer \(u\), set
\begin{equation}\label{inp:polynomials}
 P_u(\cos z)=\left(\frac{\sin(uz/2)}{u\sin(z/2)}\right)^{64},
 \qquad \mathsf P_u=P_u(1-A/64).
\end{equation}
The apparent singularities in this polynomial are removable.
The argument \(1-A/64\) is the specialization of
\(1-Q/16\) in \cite[Equation (9.11)]{B}:
for nearest neighbors \(Q=A/4\) and \(v_J^2=1/4\), so
\(A_0=Q/v_J^2=A\).  At an entry side \(m\), and at subsequent
sides \(u=mL^j\), the covariances are
\begin{equation}\label{inp:covariances}
 C_{<u}=A^{-1}(1-\mathsf P_u),\qquad
 \Gamma_u=A^{-1}(\mathsf P_u-\mathsf P_{Lu}),\qquad
 C_{\ge u}=A^+\mathsf P_u,
\end{equation}
with \(C_{<m}\) integrated at entry.  The first two inverses mean polynomial
continuation at the zero eigenvalue; the last formula acts on
nonconstant modes.  No independent-variance factor is inserted.
The same polynomials are used on free squares, with even reflection
at their walls.  Positivity, finite range, and exact local copying
apply to the polynomial accumulated and shell covariances: their
ranges are \(O(u)\) and \(O(Lu)\), respectively.  The terminal
covariance need not have finite range.  These assertions and the
derivative bounds follow from
\cite[Lemmas 8.1 and 4.1]{B} with this argument.
In particular positive-order shell and terminal estimates through
the orders used below hold uniformly when \(u\le M/D_0\), with
fixed sufficiently large \(D_0\); the plane shell diagonal is
\((2\pi)^{-1}\log L\) with an error that is a negative power of \(u\).

Use the block geometry of \cite[Lemma 4.3]{B}: with
\(p=32\), an intermediate connectivity radius is \(r_0=4p+20\)
and the polymer connectivity radius is \(\rho=r_0+8\).
These fixed radii may be enlarged for the prescribed test collars.
Shell covariance ranges are measured in the output blocks of side
\(Lu\), whereas the accumulated covariance \(C_{<u}\) has range
\(O(u)\).  The radius \(\rho\) is distinct from the tile
parameter \(R\) used in the entry construction.  Blocks have side
\(u\) on the plane and torus.  On a free square use nested translated
grids, with possible end rectangles whose side lengths lie in
\([u,2u)\).  A polymer \(X\) is a nonempty finite block set connected
by steps of sup block-index distance at most \(\rho\); \(|X|\)
counts its blocks.  Two polymers are compatible when their mutual
block-index distance exceeds \(\rho\).  The coarse closure
\(\bar X\) consists of the \(Lu\)-blocks meeting \(X\).
The geometry of \cite[Lemma 4.3]{B} supplies a fixed
\(n_0\) such that
\[
 |\bar X|\le |X|,\qquad
 |X|>n_0\ \Longrightarrow\ |\bar X|\le |X|/2;
\]
if \(|X|\le n_0\), the coarse index diameter of \(\bar X\) is
at most one.  The number of size-\(p\) polymers through a fixed
block is at most \(C_\rho^p\).

Let \(D_u(X)\) contain \(X\) and all blocks at index distance at
most two, with the fixed difference stencils and whole observation
cells included.  Activities depend on the field in \(X\) and a
fixed microscopic collar; \(D_u(X)\) is the larger testing domain.
Our activities are locally analytic functions of real fields and are
periodic under the common shift \(\psi\mapsto\psi+\omega\).
Field-even means \(K(X;-\psi)=K(X;\psi)\), a different symmetry
from the half-period parity used below.
For \(D=D_u(X)\), write \(E\) for the ambient nearest-neighbor
edge set and use the half-edge energy and boundary term
\begin{align*}
 g_D(\psi)&=\frac12\sum_{v\in D}\sum_{e\in\{\pm e_1,\pm e_2\}}
                                      |\nabla_e\psi(v)|^2,\\
 \mathcal B_u(D,\psi)&=u\sum_{\substack{\{v,w\}\in E\\|\{v,w\}\cap D|=1}}
                                             |\psi(v)-\psi(w)|^2,\\
 \mathcal T_u(D,\psi)&=\sum_{B\subset D}u^4
                     \max_{v\in B,\,|a|=2}|\nabla^a\psi(v)|^2.
\end{align*}
Physical free-wall stencils are interpreted by even extension.
Set
\[
 U_u(X,\psi)=g_D(\psi)+\delta\mathcal B_u(D,\psi)
                      +\mathcal T_u(D,\psi),\qquad
 W_u^\kappa(X,\psi)=e^{\kappa U_u(X,\psi)}.
\]
For a direction \(f\) define
\[
 \norm f_{u,D}=\max_{v\in D}
                 \{|f(v)|,u|\nabla f(v)|,u^2|\nabla^2f(v)|\}.
\]
The analytic point norm, at a real background \(\psi\), is
\begin{equation}\label{inp:point-norm}
 |F|_{u,X,\psi}=
 \sum_{r\ge0}\frac{h_0^r}{r!}
 \sup_{\norm{f_i}_{u,D_u(X)}\le1}
             |D^rF(\psi)[f_1,\ldots,f_r]|.
\end{equation}
The term with \(r=0\) is \(|F(\psi)|\).  These real derivatives
have locally convergent analytic extensions.  The full derivative
sum controls complex translations of test norm less than \(h_0\)
and obeys the product inequality.

The entry construction needs one additional, homogeneous quadratic
regulator.  Let \(B_s\) average over disjoint side-\(s\) observation
cells.  On an open union \(D\) of whole cells define
\begin{equation}\label{inp:observation-energy}
 E_D(y)=\inf_v\left\{\alpha^{-2}\langle v,A_Dv\rangle
                                     +\norm{B_sv-y}^2\right\},
\end{equation}
with free constants on every component.  Here \(A_D\) retains only
internal bonds; it differs at artificial boundaries from the
half-edge allocation \(g_D\).  Choose constants
\(1<p_1<p_2\) and \(0<h_*<1/p_2\), with the strict reserves required
by the entry estimates.  At scale \(u=mL^j\), put
\begin{equation}\label{inp:variational-regulator}
 V_j(X,\psi)=\frac12\sup_{\zeta,e}
 \left\{h_*E_D\bigl(B_s\alpha(\psi+\zeta)+e\bigr)
       -p_1^{-1}\bigl(\norm\zeta_{<u}^2+\norm e^2\bigr)\right\}.
\end{equation}
The Cameron norm \(\norm\zeta_{<u}\) belongs to the ambient
accumulated covariance \(C_{<u}\), including the pieces integrated
before entry.  Only its restriction to \(D\) enters the supremum.
The variable \(e\) is the independent unit observation noise,
included once.  For a singular covariance we use its Gaussian
Hilbert space, without inverting a null direction.

For a polymer activity \(K\) our norm is
\begin{equation}\label{inp:norm}
 \norm K_{u,\mathfrak A,V}=
 \sup_{X,\psi}
 \frac{\mathfrak A^{|X|}|K(X)|_{u,X,\psi}}
      {W_u^\kappa(X,\psi)e^{V_j(X,\psi)}}.
\end{equation}
The supremum includes all real fields.  The activity weight
\(\mathfrak A\) is unrelated to the Laplacian \(A\).  A fixed
finite collection of stronger weights, slightly larger regulator
coefficients, and analytic-radius reserves may be chosen throughout.

\begin{proposition}[Regulator input]\label{inp:regulator}
The regulator in \eqref{inp:variational-regulator} is finite,
nonnegative, homogeneous quadratic, and constant invariant.  It is
monotone under support enlargement and additive on compatible
separated supports, and
\[
 0\le V_j(X,\psi)\le CE_D(B_s\alpha\psi)
                            \le C\langle\psi,A_D\psi\rangle.
\]
The shell and terminal regulator estimates of
\cite[Lemmas 4.4 and 9.2]{B} hold with
\(W_u^\kappa\) replaced by \(W_u^\kappa e^{V_j}\), at a cost
\(C^{|X|}\).  For a small connected input the output retains the
half-exponent reserve \(W_{Lu}^{\kappa/2}\); the coefficient of the
observation regulator is unchanged.  The localizations and analytic
partition map retain their contraction and nonlinear estimates in
\eqref{inp:norm}.
\end{proposition}

The key covariance bound underlying this input is
\(A^{1/2}C_{<u}A^{1/2}\le I\).  Restricting to an open support
only lowers the minimized energy in \eqref{inp:observation-energy}.
It is therefore the homogeneous energy that is used in
\eqref{inp:variational-regulator}, also when measuring activities in
a sector with a magnetic connection.

\subsection{The exact partition map and its balanced coordinates}

At a finite block scale define
\begin{equation}\label{inp:partition-function}
 \mathcal Z_u(K;\psi)=
 \sum_{\mathcal C\ \mathrm{compatible}}\prod_{X\in\mathcal C}K(X;\psi),
\end{equation}
including the empty collection with value one.  Equivalently one sums
over block subsets and multiplies the activities of their
\(\rho\)-connected components.  A compulsory insertion activity
\(I\) uses the degree-one functional
\begin{equation}\label{inp:insertion-partition}
 \mathcal Z_u^{(1)}(K,I;\psi)=
 \sum_X I(X;\psi)
 \sum_{\substack{\mathcal C\ \mathrm{compatible}\\
                  Y\ \mathrm{compatible\ with}\ X\ (Y\in\mathcal C)}}
                                      \prod_{Y\in\mathcal C}K(Y;\psi).
\end{equation}
It is the coefficient of \(\tau\) in
\(\mathcal Z_u(K+\tau I;\psi)\), with \(\tau^2=0\).
An input polymer is called small when \(|X|\le n_0\), and large
when \(|X|>n_0\), with the geometric threshold fixed above.

\begin{proposition}[Local partition map]\label{inp:partition-map}
In the norms above, shell integration and designated local
subtractions give an exact identity
\begin{equation}\label{inp:exact-map}
 \E\mathcal Z_u(K;\psi+\zeta)
   =\prod_{b\in\mathcal B_{Lu}}(1+a_b)
                                      \mathcal Z_{Lu}(K';\psi),
 \qquad \operatorname{Cov}(\zeta)=\Gamma_u.
\end{equation}
The map is complex analytic in small activity inputs, uniformly in
finite volume, and its nonlinear remainder is bounded by
\(C_{L,\mathfrak A}\varepsilon^2\) when the activity norm and
separately measured singleton coefficients are at most \(\varepsilon\).
It preserves the common-shift period and all symmetries respected by
the localizations.  The linear map restricted to large supports has
arbitrarily small operator norm, relative to the input activity norm,
after choosing \(\mathfrak A\) sufficiently large at fixed \(L\).

For small field-even neutral bulk inputs, subtraction of the constant
and quadratic gradient term gives \(CL^{-3}\) before counting fine placements.
Subtracting only the constant gives \(CL^{-2}\) for even inputs
and \(CL^{-1}\) without evenness.  There are respectively
\(O(L^2)\), \(O(L)\), and \(O(1)\) placements at a bulk block,
a specified wall, and a fixed mark.  At \(\alpha^2=8\pi\), the
fundamental charge has gain \(O(L^{-2})\); for field-even bulk inputs,
subtraction of its constant-background cosine coefficient gives an
additional \(L^{-1}\).
Higher charges have summable stronger gains.
\end{proposition}

These are Proposition 4.8, Lemma 4.5, Corollary 4.7, and
Proposition 8.3 of \cite{B}.
Before localization, Gaussian integration groups the union of the
coarse closures into components connected at radius \(r_0\).
Different such components are beyond the shell covariance range.
The subtracted pieces are assigned to coarse singletons with their
original exclusion information retained.  Only designated
field-independent coefficients enter \(a_b\).  In particular, a
translated quadratic polynomial may retain its nonzero value at zero
field inside the activity.  The same algebra applies in degree one
in \(\tau\); no smallness of the vector \(I\) is required to apply
the derivative of the map at a small background.  Coefficients copy
from the plane whenever their predecessor neighborhoods fit inside
the domain with the fixed clearance supplied by the covariance range.

For a full bulk block define the two singleton functions
\[
 e_B^0(\psi)=\frac12\sum_{v\in B}
              \sum_{e\in\{\pm e_1,\pm e_2\}}|\nabla_e\psi(v)|^2,
 \qquad c_B^\alpha(\psi)=u^{-2}\sum_{v\in B}\cos(\alpha\psi(v)).
\]
Their coefficients in a general activity are denoted by \(t/2\)
and \(z\): the singleton part is \(t e_B^0/2+zc_B^\alpha\).
Denote the remaining activity by \(\mathcal R\).  The linear map is
triangular,
\[
 (t,z,\mathcal R)_+
   =(t+P_j^t\mathcal R,\lambda_jz+P_j^z\mathcal R,C_j\mathcal R)
       +O\bigl((|t|+|z|+\norm{\mathcal R})^2\bigr),
 \qquad \norm{C_j}\le\theta<1,
\]
where \(\lambda_j=L^2e^{-\alpha^2\Gamma_u(0,0)/2}\).
At marginality the product \(\ell_j=\prod_{h\ge j}\lambda_h\)
converges.  Set \(C_{h:j}=C_{h-1}\cdots C_j\) for \(h>j\)
and \(C_{j:j}=I\).  The balanced coordinates are
\begin{align}\label{inp:balance}
 T&=t+\sum_{h\ge j}P_h^tC_{h:j}\mathcal R,\nonumber\\
 Z&=\ell_jz+\sum_{h\ge j}\ell_{h+1}P_h^zC_{h:j}\mathcal R.
\end{align}
Both projection series have geometric tails.  Their linear map is
\((T,Z,\mathcal R)\mapsto(T,Z,C_j\mathcal R)\).

The exact Gaussian curves of
\cite[Lemma 10.3]{B}
are obtained by initializing the interaction
\[
                         \exp\left(\frac d2\sum_Be_B^0\right)
\]
without a lattice factor.  For each fixed entry scale they have
\(T_j^g(d)=d+O(d^2)\), \(Z_j^g(d)=0\), and
\(\mathcal R_j^g(d)=O(d^2)\), analytically on a disk independent
of the step.  In the following coordinate formulas \(u\) denotes
the Gaussian scalar coordinate rather than a physical block side.
Define
\[
 u=(T_j^g)^{-1}(T),\qquad
 B=\mathcal R-\mathcal R_j^g(u).
\]
Thus \((u,0,0)\) is an exact fixed-coordinate Gaussian trajectory.
Split \(B=B_e+B_o\) into its even and odd parts under the
half-period shift \(\psi\mapsto\psi+\pi/\alpha\).

\begin{proposition}[Balanced parity map]\label{inp:balanced-map}
For the critical reference and small, even, translation- and
square-covariant plane activities, the preceding coordinates satisfy
\begin{align}\label{inp:parity-equations}
 u_+-u={}&-H_jZ^2+
 O\bigl((|u|+Z^2)(Z^2+\norm{B_e})+\norm{B_e}^2
                    +\norm{B_o}(|Z|+\norm{B_o})\bigr),\nonumber\\
 Z_+-Z={}&-b_juZ+
 O\bigl((u^2+Z^2+\norm{B_e})|Z|
       +( |u|+Z^2+\norm{B_e}+\norm{B_o}^2)\norm{B_o}\bigr),\nonumber\\
 \norm{B_{e,+}}\le{}&\theta_1\norm{B_e}
                              +C(|Z|+\norm{B_o})^2,\nonumber\\
 \norm{B_{o,+}}\le{}&\theta_1\norm{B_o}
                       +C(|u|+Z^2+\norm{B_e})|Z|,
 \qquad \theta_1<1.
\end{align}
The estimates include their differentiated analytic Taylor bounds.
The coefficients are bounded above and away from zero, and
\begin{equation}\label{inp:coefficient-limits}
                         H_j\to H>0,\qquad b_j\to2\log L.
\end{equation}
These are the map in Lemma 10.6 and the coefficient
calculation in Proposition 8.5 of \cite{B}.
\end{proposition}

This proposition is a local analytic input, conditional on entry into
and continued residence in its small tube.  The uniformity needed
when the entry scale varies, the absence of exit, and the quantitative
rate in \eqref{inp:coefficient-limits} are established in the bulk
argument below.

Finally, we record the source estimates used with these local maps.
For the one-center connection \(l\) of Lemma~\ref{inp:duality},
let \(r\) minimize \(\norm{\partial r+l}^2\) on the free dual
square and put \(\eta_l=\partial r+l\).  If \(G_n\) is the
inverse of the killed unit-edge Laplacian on the interior spin
vertices, then orthogonal energy minimization gives
\begin{equation}\label{inp:connection-energy}
                         \norm{\eta_l}^2=(2\pi)^2G_n(0,0).
\end{equation}
By \cite[Lemmas 4.2 and 5.2]{B},
\begin{equation}\label{inp:source-lifts}
 |\nabla^a\eta_l(v)|\le C_a(1+\operatorname{dist}(v,0))^{-1-a},
                         \qquad 0\le a\le5.
\end{equation}
On a source-free enlarged rectangle the connection is an exact real
gradient; the assertion includes the reflected stencils at a free
wall.  Plane and finite-square connection gradients have the same
local limit.  The stronger finite-volume difference estimate needed
for the magnetic calculation follows there by subtracting the common
finite-range pieces and summing the positive-order tails.

All norm choices follow a fixed order: geometry and test collars;
the observation-regulator reserves;
\(\delta\), then \(\kappa\) with the necessary integration reserves;
the analytic radius \(h_0\); a sufficiently large power-of-two
block factor \(L\); activity weights; and the small activity tube.
For the XY entry construction, the observation-regulator choices must
also satisfy the stronger reserve condition in Proposition~\ref{prop:entry}.
This condition is imposed before fixing \(L\) and the weights.
The subsequent tile parameters and observation-cell size are chosen
after these map parameters.  Constants may depend on parameters
already fixed, but never on the volume or the running scale unless
explicitly indicated.
